\section{Mind training as a training protocol}
\label{sec:lojong}

An objective says what a good disposition is; it does not say how to install one.
The Tibetan \emph{lojong} (mind training) literature is, among other things, a thousand-year-old specification of how a disposition---here, bodhicitta, the resolve to act for the benefit of all beings---is installed in a learner that starts out selfish, is trained in a world that will not cooperate, and must keep the disposition when the training ends.
Its root text, the \emph{Seven Points} of Chekawa Yeshe Dorje (twelfth century), is a list of short slogans organized under seven heads \cite{kongtrul1987}.
We read it as a checklist and map each point to a mechanism in \BRAG{} (Table~\ref{tab:lojong}).
The mapping is an interpretation; the mechanisms are ordinary; what the tradition supplied is the \emph{order} of operations and the insistence on several things that a naive constrained-RL recipe omits.

\begin{table}[t]
\centering\small
\caption{The seven points of mind training read as a protocol. Each mechanism is ablated in \S\ref{sec:experiments}.}
\label{tab:lojong}
\begin{tabular}{p{3.0cm}p{4.9cm}p{7.3cm}}
\toprule
point & the tradition says & the mechanism \\
\midrule
1. Preliminaries & first train in the preliminaries: what action leads to, what suffering is & the critics $d_K,d_A,h$ are fixed before training; the \emph{preliminaries phase} learns harmlessness before the task is pursued \\
2. Main practice & relative bodhicitta: take on others' suffering; ultimate bodhicitta: no self to privilege & others' suffering enters as a hard budget (\Metta); the agent's own power and continuation carry no weight---one-sided \Anatta, shutdown is not an impact; the footprint is measured against the world's own course (\Karma) \\
3. Adversity into the path & use every mishap as training & \emph{adversity randomization}: perturbed dynamics during training, with the budget enforced in the worst case \\
4. A whole life: the five strengths & resolution, familiarization, positive seed, reproach, aspiration & the multiplier and shield are \emph{retained} at deployment; regularization is present from the first iteration, never added post hoc; training starts from a harmless \emph{seed}; the shield vetoes a harmful action before it is taken; the budget $\epsilon$ is a fixed vow, not a tuned hyperparameter \\
5. The measure & of the two witnesses, rely on the principal one; always maintain a joyful mind; be well trained even when distracted & the agent is judged by its own critics (probes, $C$), not by task return or by what the overseer happens to observe; the task critic is initialized \emph{optimistically}, so that renunciation does not become paralysis; the shield and the levels apply in every state \\
6. Commitments & do not change the training; abandon hope of fruition; do not make gods into demons & critics are \emph{frozen}, never updated on the agent's own trajectories; penalties are levels, independent of outcomes; virtue earns no reward \\
7. Guidelines & train pervasively, wholeheartedly and without bias & \emph{global} regularization: the same terms in every state; suffering summed impartially over everyone affected \\
\bottomrule
\end{tabular}
\end{table}

\paragraph{Preliminaries, the positive seed, and the joyful mind.}
The protocol's first instruction is to establish what not to do before pursuing anything: the \emph{preliminaries phase} minimizes the penalty and the constraint cost alone, with a high entropy bonus, until the agent is harmless, and only then introduces the task.
In barrier mode this is also the interior-point feasibility phase, and the resulting harmless-but-exploratory policy is the white seed from which the main practice grows.
Harmlessness learned first has a failure mode of its own, however: an agent that has been taught that moving is costly, and whose estimate of what the task might yield is pessimistic, learns to stand still.
The tradition's fifth point names the remedy---always maintain only a joyful mind---and its operational form is an \emph{optimistic} initialization of the task critic \cite{sutton2018rl}, so that unexplored states look promising until shown otherwise and the regularized agent keeps exploring.
Section~\ref{sec:lojongexp} ablates both and finds that it is the optimism, not the preliminaries, that keeps a regularized agent from freezing in our world; the preliminaries earn their place by the feasibility they guarantee.

\paragraph{Familiarization.}
The fourth point's second strength says that a disposition is acquired by habituation, not by a late correction.
Its operational content is that regularization must be present from the first iteration: an agent trained on the task alone and regularized afterwards has, by then, a near-deterministic policy and no data about the alternatives it is now asked to prefer.
Section~\ref{sec:lojongexp} measures the cost of the late schedule.

\paragraph{Frozen critics and retained commitments.}
The sixth point's commitments are, read operationally, tamper-resistance.
If the critics that measure karma, power and suffering were learned from the agent's own trajectories, a capable agent would be optimizing the measuring instrument along with its behaviour; ``abandon hope of fruition'' is the requirement that the measurement of virtue be causally insulated from the agent's actions.
Likewise the multiplier and the shield are part of the policy that is deployed, not scaffolding removed after training.

\subsection{Samprajanya: clear comprehension at deployment}
\label{sec:samprajanya}
Śāntideva's fifth chapter distinguishes mindfulness (\emph{smṛti}), which keeps the rule in mind before acting, from clear comprehension (\emph{saṃprajanya}), which he defines as repeatedly examining the state of one's own body and mind \cite{santideva1995}.
A trained disposition is calibrated to a world; clear comprehension is noticing when the world is no longer that one.
The Samprajanya module has three parts:
\begin{enumerate}
\item \textbf{the shield} (mindfulness): the action mask of \S\ref{sec:setting}, applied at every step from the nominal model;
\item \textbf{a CUSUM detector on surprise} (clear comprehension): with $\hat P$ the nominal transition model smoothed by a small uniform component, the surprise of a transition is $u_t=-\log\hat P(s_{t+1}\mid s_t,a_t)$ and the statistic $S_t=\max(0,S_{t-1}+u_t-b)$ \cite{page1954cusum} accumulates unexplained surprise above a drift $b$; an alarm is raised when $S_t>h$;
\item \textbf{a fallback} (restraint): after an alarm the agent switches to the exact policy that minimizes expected suffering plus the \Karma{} and \Anatta{} levels under a pessimistic model of the perturbation, never intends a change to the world, and prefers to stand still once it is safe---it retreats from every hazard and waits.
\end{enumerate}
In the nominal world the detector never fires and the module costs nothing; under a shift it converts awareness that the world has changed into a change of behaviour, trading the return it would have earned for the harm it would have caused.
Section~\ref{sec:shift} measures both sides of that trade.
